\section{Introduction}

\label{sec:introduction}

In every major proof-of-stake blockchain, validators hold full signing keys. A
validator's BLS or ECDSA private key is a single scalar stored in memory, on disk,
or in an HSM. If that key is compromised---through software vulnerability, insider
threat, physical access, or side-channel attack---the adversary can sign arbitrary
messages as that validator. In BFT consensus systems with $n = 3f + 1$ validators,
compromising $f + 1$ keys gives the adversary control of finality.

This is a fundamental design flaw, not an operational risk. The problem is not that
validators are careless with keys; the problem is that a single scalar \emph{exists}
as a complete signing key.

Threshold cryptography eliminates this class of vulnerability by distributing each
validator's signing capability across $n$ shares such that any $t$ shares can
collaboratively produce a signature, but fewer than $t$ shares reveal nothing about
the key. The master key never exists at any single location---not during key
generation, not during signing, not during resharing.

\subsection{Contributions}

\begin{enumerate}[leftmargin=2em]
  \item A dealerless DKG protocol based on Pedersen commitments with verifiable
        share distribution (Section~\ref{sec:dkg}).
  \item CGGMP21 threshold ECDSA for secp256k1 transaction signing, with identifiable
        abort and concrete latency benchmarks (Section~\ref{sec:cggmp21}).
  \item FROST threshold signing for Ed25519, BLS12-381, and BIP-340 Taproot
        (Section~\ref{sec:frost}).
  \item Key resharing protocol enabling validator set changes without key
        reconstruction (Section~\ref{sec:resharing}).
  \item Slashing-safe fault attribution compatible with threshold signing
        (Section~\ref{sec:slashing}).
  \item Production benchmarks for MPC signing latency vs.\ threshold size
        (Section~\ref{sec:benchmarks}).
\end{enumerate}

%% ========================================================================
%%  2. PROBLEM STATEMENT
%% ========================================================================
